In recent years, network models have become a widespread method for studying the structure of multivariate relationships within psychological systems. The goal of these models is generally to characterize the direct associations among different variables that are theorized to be interrelated with regard to some psychological process or construct. These models have been applied within both cross-sectional and longitudinal contexts, where in the former the goal is often to describe which variables have direct relationships across multiple individuals, and in the latter the goal is often to describe how different variables relate with one another within a single individual over time. Multilevel implementations of longitudinal models—termed temporal networks—have also been employed in the study of whether those within-person processes hold across groups of individuals.

While these network models have been shown to have a variety of diverse applications, they are limited by the fact that they only consider pairwise relationships among sets of variables. Specifically, they don’t take into consideration more complex relational structures such as those characterized by moderation effects. Moderation analysis is a common procedure within psychological research, and can be used to help elucidate the contexts within which different relationships may emerge, or be observed, between different variables. For instance, we may be interested in whether the relationship between some variable $X$ and some variable $Y$ is dependent on the value of some third variable $Z$. It is well-known that moderator variables can be crucial in understanding how different variables relate with one another, and that leaving out such variables can lead to a distorted understanding of the underlying process or causal structure. Thus, my goal for this research is to extend the psychological network framework to include moderator variables, as well as provide statistical tools and software to facilitate testing such models using psychological data. 

The focus of the present research is particularly on moderated temporal networks—that is, testing moderator variables within network models applied to longitudinal and time series data. It is worth noting that the methods described herein can be applied to the analysis of cross-sectional data as well, however the focus will be on data collection for repeated observations of a single individual over time, as well as how models within this context can be extended to datasets with multiple subjects. I will begin by describing some more background on psychological network models and moderation analysis, then discuss some of the basic time series models that are used in specifying temporal networks, as well as how these methods can be combined to create moderated temporal networks. Finally I will conclude by presenting some of the extensions and considerations that I investigate in an R package designed to implement such models and allow researchers to test them with their own data. 

\apasection{Psychological Networks}
Network models have gained great popularity in psychological research. The idea behind networks is to describe the structure of complex systems using concepts from graph theory, which is a field of study that emerged to characterize information processing in computer systems and information exchange in transportation systems. The core concepts of graph theory are captured by the notion of components, or ‘nodes’, whose relationships are depicted as links, or ‘edges’, that denote the presence or absence, and possibly the strength and direction, of those relationships. The utility of network models comes from their capacity to represent dependencies among variables that, together, are assumed to constitute the operation of a complex system. Given the complexities inherent in psychological processes, network models offer a new set of tools for studying them.

Psychological networks are models that provide both a statistical and theoretical representation of a complex psychological system or construct. In these models, there are a set of `nodes` which represent psychological variables (such as mood states, attitudes, symptoms, etc.), and then a set of `edges` which represent statistical relationships among those variables. One domain in which these models have recently become popular is in the \textit{Network Approach to Psychopathology} [CITE], wherein researchers are conceptualizing and studying mental disorders as networks of directly interacting and mutually reinforcing symptoms. From this perspective, symptoms of mental disorders such as Major Depression (MD) co-occur because of their direct causal relationships rather than being products of a latent, common cause. For example, instead of treating symptoms such as insomnia, fatigue, and dysphoria as products of an underlying neurochemical imbalance or brain disorder, the network approach to psychopathology considers the direct relationships of insomnia $\rightarrow$ fatigue $\rightarrow$ dysphoria as a contributor to the experience of MD. 

The origin of these models from a mathematical standpoint comes from graph theory, which emerged in the XXXs with the goal of characterizing transportation networks (such as railroads) and computer systems. Further statistical development of these methods comes from Probabilistic Graphical Models [lauritzen], wherein connections between nodes are represented as probabilistic, weighted associations---particularly as \textit{partial correlations}. The most widely-used model of this type is known as the \textit{Gaussian Graphical Model} (GGM), which describes partial correlations among a set of multivariate-normal continuous variables. Although these models don't directly encode causal relations between variables, the use of partial correlations is aimed at uncovering unique pairwise relationships that, in exploratory contexts, may be useful in assessing possible causal pathways. 

Another approach to network models used by researchers are Bayesian Networks, which are a type of Directed Acyclic Graph (DAG) that directly encode causal relationships. While these models are not often used within the context of psychological networks, directed networks are often still estimated within longitudinal or repeated-measure studies. Directed and undirected networks have distinct properties which will be discussed below. 

\apasubsection{Types of Networks}
The first distinction between different types of network models is between \textit{undirected} and \textit{directed} networks. The difference between these types of networks lies in whether or not the direction of the association between a given pair of variables is encoded in the edge connecting them.

\apalevelthree{Undirected Networks}
The most basic form of network models used in psychology is referred to as the Gaussian Graphical Model (GGM). This model is used to describe relationships among a set of continuous, normally distributed variables. In GGMs, each variable is represented as a node in the network, and their relationships (edges) are typically represented as partial correlations. Specifically, the GGM is defined as the standardized inverse covariance matrix of the set of variables—i.e., the partial correlation matrix. The model is thus defined as: [EQUATION]. Where this equal that, etc.

There are two primary reasons for using the partial correlation matrix, rather than the zero-order correlation matrix, to define the GGM: First, the goal is to characterize the conditional dependencies among the variables under consideration. That is, we are interested in understanding the unique relationships among any given pair of variables, controlling for the relations they have with other variables in the system. Importantly, this means that we are also interested in whether certain pairs of variables are conditionally independent of one another, given their other multivariate relationships. This is arguably one of the primary objectives of network modeling, in that we are often interested in revealing the core structure underlying the system at hand. Thus, partial correlations offer a distinct advantage in this sense, given that their interpretation affords an understanding of these conditionally dependent relationships. 

Second, and corresponding to the first point, partial correlations are typically smaller than their zero-order counterparts, and may therefore lead to a more sparse model than if zero-order estimates were used. This is directly related to the idea of exploring whether certain variables are conditionally independent of one another. In light of this goal, psychological network models, particularly GGMs, are often estimated with regularization techniques that involve penalized loss functions, such as the LASSO. This technique has important implications for the interpretation and estimation of network models, and will be discussed further in forthcoming work.

\apalevelthree{Directed Networks}
Another general type of network models, used specifically with time series data, is the temporal network, which is characterized by a network of directed relationships, typically defined as regression coefficients rather than undirected partial correlations. Temporal networks are much more complicated to specify than undirected networks of cross-sectional data, and they also don’t have as straightforward of an interpretation. 

From graph theory, the standard model used to depict a temporal network would be a Directed Acyclic Graph (DAG), which is a network that consists of nodes with arrows representing directed, causal relationships between them. A critical aspect of these networks, however, is that the graph only preserves information about the conditional dependence structure of the constituent variables if there are no self-loops in the model; i.e., no autoregressive components. This assumption is often not plausible for psychological data, however, as repeated measurements of a single variable often contain unique predictive information about subsequent measurements. Importantly, this means that true DAGs are not often used to construct temporal networks with psychological data, and instead these network models may be more properly understand as visual representations of multivariate multiple regression models. 

The key implication of the fact that most temporal networks in psychology are not DAGs is that most temporal networks in psychology do not function as true graph-theoretic models. However, there are certain undirected graphical models that can be estimated in conjunction with temporal network models, and these will be discussed further in this paper. Specifically, we can use the theory behind GGMs to construct contemporaneous networks from repeated measurement models, and in cases where we have repeated measurements taken across multiple individuals, between-subject networks can also be constructed. More details about these models will be discussed in later sections of this paper.

\apasubsection{Interpreting Networks}
Importantly, network models encode the conditional (in)dependence structure among a given set of observed variables. We use networks to investigate whether specific variables have unique relationships with one another after considering their relationships with some other variables of interest.

\apasubsection{Extending Network Models}
One important aspect of network models is that they only encode information about pairwise relationships among variables, and do not take into account how some relationships may vary in accordance with changes in other variables. An implicit assumption of these models is that pairwise relationships between variables will be the same across all values of both variables. In temporal networks, this means that we assume, for instance, that variable X has the same predictive relationship with subsequent values of variable Y across the entire observation window. Although in some circumstances this may not be an entirely unreasonable assumption, it seems particularly problematic in many cases where the relationships among variables may change in accordance with changes in, say, situational stress or other environmental factors. Thus, an important goal for extending the framework of network modeling, and particularly temporal networks, is to afford researchers the ability to test complex hypotheses---such as that of moderation or interaction effects---as well as use exploratory techniques to investigate multivariate causal structures in psychological data. 

Next, I will outline the basic idea behind moderation as well as motivate reasons why it is important to investigate these types of relationships in psychological networks. 

It is an important goal for psychologists to develop and improve upon methods for describing complex relationships among observed variables. The dominant approach toward this end has long been \textit{latent variable modeling}, where the objective is to connect a set of observations with some unobserved latent construct(s) using methods such as confirmatory factor analysis and structural equation modeling (SEM) [CITE]. In general, these methods aim at highlighting the shared variance among variables while downplaying their conditional relationships. The \textit{network perspective} in psychology, however, offers a powerful counterpoint to this approach by framing psychological constructs as systems of interacting variables rather than manifestations of unobserved factors. Thus, while latent variable models are designed to sift through observed variation in search of underlying commonalities, network models focus on the direct associations between variables in an effort to narrow down their unique causal structure. For instance, rather than conceptualizing cognitive abilities, behavioral dispositions, and mood disorders as indicators of underlying common causes---such as general intelligence, personality traits, and psychopathology---psychological network models frame these phenomena as emergent characteristics of networks of interrelated components. 

The primary goal behind the network perspective is to characterize the structure of psychological and behavioral systems using concepts from graph theory, a field of study that emerged in the XXXs to characterize information processing in computer networks and transportation systems. Network models provide both statistical and theoretical representations of psychological constructs, wherein a set of variables (such as mood states, attitudes, symptoms, etc.) can be represented as components, or 'nodes', whose unique associations are depicted as links, or 'edges'. One domain in which these models have been applied is the \textit{network approach to psychopathology}, where mental disorders are viewed as arising from networks of directly interacting and mutually reinforcing symptoms [CITE]. From this perspective, symptoms of mental disorders such as Major Depression (MD) co-occur because of their direct causal relationships, rather than their connections to some underlying latent construct. For example, instead of treating symptoms such as insomnia, fatigue, and dysphoria as consequences of a neurochemical imbalance or brain disorder, the network approach to psychopathology considers causal pathways such as insomnia $\rightarrow$ fatigue $\rightarrow$ dysphoria to be key contributors to the experience of MD. 

%%%%%%%%%%%%%%%%%%%%%%%%%%
\apasection{Network Models in Psychology}
Network models have been emerging as powerful tools for conceptualizing and analyzing psychological processes in terms of complex systems. These types of models are used to map dynamic links among observed variables with an emphasis on the \textit{structure} of those relationships. Networks have been used to study the converging influences of complex cognitive, biological, and sociological components of psychological states, memory processes, and emotion dynamics (to name a few) within a common statistical and operational framework. For instance, network models have been used to [CITE], [CITE], [CITE], and [CITE]. Many methods for analyzing network data are still being tested and developed, and canonical graph-theoretic approaches to studying networks do not apply in a straightforward way to the domain of psychological networks. Thus, it is an important goal to further understand and expand methods that are associated with this framework and which connect it to more common statistical approaches used in psychology. 

A major area of research that has utilized network models within psychology has been focused on conceptualizing psychopathology as a network of mutually interacting symptoms, wherein the structure of their relationships are theorized to govern the experience and character of many psychological disorders. These phenomena are conceived of as being the emergent consequences of direct, causal associations (`edges') among psychological, biological, and sociological variables (`nodes'), modeled as a probabilistic network. These models are also being applied within the study of personality, wherein personality and individual differences are studied in terms of complex interactions among cognition and behavior. [EXAMPLE OF PERSONALITY RESEARCH]. These two areas of study have not seen a great deal of overlap, however, given the limited research and available statistical approaches. 

Experience sampling approaches and access to intensive longitudinal data have been important components of the popularity of these types of models. This has opened up possibilities of applying time series methods and investigating both linear and nonlinear relationships among day-to-day experiences such as mood and vulnerability to depression. Experience sampling methods generally involve having participants complete surveys and complete psychological measures throughout their daily life---this may occur only once or twice per day (as in the traditional `daily-diary' study), or anywhere from 5 to 10 times per day [CITE]. These time-intensive data are available in much greater volume now given the widespread accessibility of mobile phones and other `smart' devices that pervade our daily lives. These methods appear frequently in the psychological networks literature, and provide a potentially powerful view into intraindividual dynamics as well as those that can be observed between-subjects. 

\apasubsection{Types of Networks}
Networks are graphical models that consist of \textit{nodes} (variables) and \textit{edges} (relationships between variables). These models have their origins in graph theory from discrete mathematics and computer science, and are applied to such diverse phenomena as semantic structure, transportation systems, electrical grids, and gene-regulatory processes. There are a variety of different types of network models, although there are a few notable categories that can be distinguished at the outset. In general, there are two primary dimensions across which networks can vary: (a) Undirected versus directed, and (b) Weighted versus unweighted. 

\apalevelthree{Undirected networks}
A network is a `graph', in technical terms, 

\apalevelthree{Directed networks}
These are when time is involved.

\apalevelthree{Directedness}
Whether a network is \textit{directed} or \textit{undirected} refers to the difference between networks where time is explicitly represented (in directed networks) or is not (in undirected networks). In directed networks, the edges between nodes have arrows at the ends to indicate which variable predicts or activates the other, and thus have a temporal meaning. In undirected networks, however, the edges are simply straight lines between variables, indicating that there is believed to be some sort of relationship between them, but without reference to a particular causal hypothesis about which variable influences the other or whether they have a reciprocal connection.

\apalevelthree{Weightedness}
Whether the edges of a network are \textit{weighted} or \textit{unweighted} is independent of its directedness, as it instead refers to whether the edge is being used to indicate some quantitative relationship between variables (e.g., similarity, correlation), or a qualitative one (e.g., whether the two components are connected or not). Different application may call for different types of networks---for instance, semantic and phonological network models are typically undirected and unweighted, networks for psychological disorders may frequently be directed and weighted. For the remainder of this paper, I will only be discussing \textit{weighted} networks, but will also cover some of their undirected and directed counterparts.

Weighted networks often use some type of similarity or distance metric to denote the meaning of their edges. Broadly speaking, models that use some similarity-, correlation-, or regression-based definition of an edge have been referred to as \textit{probabilistic graphical models} \cite{lauritzen1996graphical}. These models are called `probabilistic' because the edges define the some statistical quantity that describes the strength (and possibly direction) of relation between any given pair of nodes. There are a number of different ways that a researcher could construct such a network [GIVE EXAMPLES HERE]. In the following sections, I will describe in greater detail some of the more widely-used approaches, particularly within the domain of clinical and personality psychology.